% !TeX root = ../main.tex
% ============================================================
%  sections/04c_follower.tex
% ============================================================

%-----------------------------------------------------------------
\subsection{Supply Chain Operator's Problem}
\label{sec:model:follower}

Given the tax index $\taxindex\in\mathcal{T}^{\tau}$ announced by the
regulator, the supply chain operator solves a mixed-integer programme
whose feasible region $\setX$ is defined by constraints (C1)--(C10) in
\Cref{sec:model:c1,sec:model:c2,sec:model:c3,sec:model:c4,%
sec:model:c5,sec:model:c6,sec:model:c7,sec:model:c8,sec:model:c9,%
sec:model:c10}. The objective, detailed in \Cref{sec:model:follower_obj},
combines total cost and carbon-tax payment.

%- - - - - - - - - - - - - - - - - - - - - - - - - - - - - - - -
\subsubsection{Objective function}
\label{sec:model:follower_obj}

Given $\taxindex$, the operator minimises the sum of its total cost
$\Cost(\bz)$ and carbon-tax payment:
\begin{equation}
\textup{(FL)}\quad
\min_{\bz\in\setX}\;
\Cost(\bz) + \gridstep\,\taxindex\,\Emis(\bz),
\label{eq:follower}
\end{equation}
with optimal-response correspondence
\begin{equation}
\Sol{\taxindex} \defeq
\argmin_{\bz\in\setX}\;\Bigl\{\Cost(\bz) +
\gridstep\,\taxindex\,\Emis(\bz)\Bigr\}.
\label{eq:reaction_set}
\end{equation}

The operator's total cost is given by
\begin{align}
\Cost(\bz)
={}&
  \sum_{i\in\setF\cup\setW} c^{\mathrm{area}}_i\,YC_i
  +\sum_{f\in\setF}\Bigl(
     \sum_{g\in\setG} C^{\mathrm{p}}_g\,\lambda^{\mathrm{p}}_{fg}
    +\sum_{\ell\in\setL} C^{\mathrm{r}}_{\ell}\,\lambda^{\mathrm{r}}_{f\ell}
  \Bigr)
\nonumber\\
&+
  \sum_{i\in\setN}\sum_{m\in\setMrd}
     F^{\mathrm{m}}_m\,K_{im}
   +\sum_{h\in\setH}\sum_{t\in\setT}
     F_h\,y_h
\nonumber\\
&+
  \sum_{(p,g)\in\Omega^{\mathrm{p}}}\sum_{f,t}
     c^{\mathrm{p}}_g\,q^{\mathrm{p}}_{pfgt}
   +\sum_{(p,\ell)\in\Omega^{\mathrm{r}}}\sum_{f,t}
     c^{\mathrm{r}}_{\ell}\,q^{\mathrm{r}}_{pf\ell t}
\nonumber\\
&+
  \sum_{p\in\setP}\sum_{i\in\setNv}\sum_{t}
     c^{\mathrm{h}}_p\,s^{\mathrm{inv}}_{pit}
\nonumber\\
&+
  \sum_{u\in\setU}\sum_{s\in\setS}\sum_{t}
     c^{\mathrm{u}}_{su}\,X^{\mathrm{out}}_{ust}
   +\sum_{k\in\setK}\sum_{r\in\setR}\sum_{t}
     c^{\mathrm{k}}_k\,X^{\mathrm{out}}_{krt}
\nonumber\\
&+
  \sum_{(i,j)\in\setA}\sum_{m\in\setM_{ij}\cap\setMrd}\sum_{t}
     c^{\mathrm{veh}}_m\,r_m\,\delta_{ij}\,Q_{ijmt}
   +\sum_{(q,i,j)\in\Phi}\sum_{m\in\setM_{ij}\setminus\setMrd}\sum_{t}
     c^{\mathrm{tr}}_m\,\delta_{ij}\,\omega_q\,x_{qijmt}
\nonumber\\
&+
  \sum_{(q,i,j)\in\Phi}\sum_{m\in\setM_{ij}}\sum_{t}
     c^{\mathrm{hub}}_m
     \bigl(\xi^{\mathrm{o}}_{ij}+\xi^{\mathrm{d}}_{ij}\bigr)
     \omega_q\,x_{qijmt},
\label{eq:cost}
\end{align}

The first term is the fixed cost of selecting a supplier, plus the
area-proportional installation cost $c^{\mathrm{area}}_i$ charged on the
floor area $YC_i$ installed at an opened factory or warehouse -- there
is no separate fixed cost for opening a factory or warehouse, since
$YC_i\ge\underline{A}_i\,y_i$ already forces a nonzero installation cost
whenever $y_i=1$ (\Cref{eq:area_bounds}) -- plus the cost of installing
a production or remanufacturing technology at each factory. The second term is the road
fleet investment cost and the hub access fee, the latter charged every
period a hub is selected. The
third term is the variable cost of production and remanufacturing, and
the fourth is the inventory-holding cost on finished-product stock at
warehouses -- factories and recovered items are never held in
inventory (\Cref{sec:model:c1}), so neither carries a holding-cost
term. The fifth term is the acquisition cost of
raw materials and recovered items. The sixth term is transport: road
transport is charged per trip, scaled by the return factor $r_m$ for the
empty return leg, while outsourced transport is charged per
kilogram-kilometre shipped. The seventh and final term is the
hub-handling charge, applied once for each hub endpoint of an arc.

Total emissions are given by
\begin{align}
\Emis(\bz)
={}&
  \sum_{i\in\setF\cup\setW} e^{\mathrm{o,f}}_i\,YC_i
\nonumber\\
&+
  \sum_{(p,g)\in\Omega^{\mathrm{p}}}\sum_{f,t}
     e^{\mathrm{p}}_g\,\omega_p\,q^{\mathrm{p}}_{pfgt}
   +\sum_{(p,\ell)\in\Omega^{\mathrm{r}}}\sum_{f,t}
     e^{\mathrm{r}}_{\ell}\,\omega_p\,q^{\mathrm{r}}_{pf\ell t}
\nonumber\\
&+
  \sum_{u\in\setU}\sum_{s\in\setS}\sum_{t}
     e^{\mathrm{u}}_{su}\,\omega_u\,X^{\mathrm{out}}_{ust}
\nonumber\\
&+
  \sum_{(q,i,j)\in\Phi}\sum_{m\in\setM_{ij}}\sum_{t}
     e^{\mathrm{tr}}_m\,r_m\,\omega_q\,\delta_{ij}\,x_{qijmt}.
\label{eq:emissions}
\end{align}

The first term is the opening emission: an amount proportional to the
installed area $YC_i$ for an opened factory or warehouse -- suppliers
and hubs carry no opening emission of their own; a supplier's emission
footprint is captured entirely by the raw-material acquisition term
below. The second term is the production and
remanufacturing emission, per kilogram of output.
The third term is the raw-material acquisition
emission, per kilogram shipped. The fourth and final term is transport,
scaled by the same return factor $r_m$ used in the transport cost
(\Cref{eq:cost}): the empty return leg of a road delivery still burns
fuel and emits CO$_2$, even though it carries no cargo; $r_m=1$ for
outsourced modes, which have no return leg attributable to the
follower.

%- - - - - - - - - - - - - - - - - - - - - - - - - - - - - - - -
\subsubsection{(C1)~~Flow conservation and inventory balance}
\label{sec:model:c1}

These constraints enforce material balance at every node type in every
period. Finished-product inventory follows a standard stock-balance
recurrence at warehouses -- factories carry no inventory of their own
and operate as pure flow-through nodes for finished products, mirroring
the hub cross-docking pattern below; recovered items are likewise never
buffered, at warehouses or anywhere else, and pass through warehouses
as pure flow-through too; markets satisfy demand exactly with no
on-site inventory; and hub terminals operate as pure cross-docking
nodes:
\begin{subequations}\label{eq:c1_balance}
\begin{alignat}{2}
X^{\mathrm{in}}_{pft}
    + \sum_{g:(p,g)\in\Omega^{\mathrm{p}}} q^{\mathrm{p}}_{pfgt}
    + \sum_{\ell:(p,\ell)\in\Omega^{\mathrm{r}}} q^{\mathrm{r}}_{pf\ell t}
  &= X^{\mathrm{out}}_{pft},
&& \forall p\in\setP,\,f\in\setF,\,t\in\setT,
\label{eq:bal_factory}\\
s^{\mathrm{inv}}_{pwt}
  &= s^{\mathrm{inv}}_{pw,t-1}
    + X^{\mathrm{in}}_{pwt} - X^{\mathrm{out}}_{pwt},
&& \forall p\in\setP,\,w\in\setW,\,t\in\setT,
\label{eq:bal_warehouse}\\
X^{\mathrm{in}}_{kwt}
  &= X^{\mathrm{out}}_{kwt},
&& \forall k\in\setK,\,w\in\setW,\,t\in\setT,
\label{eq:bal_warehouse_rp}\\
X^{\mathrm{in}}_{prt}
  &= d_{rpt},
&& \forall p\in\setP,\,r\in\setR,\,t\in\setT,
\label{eq:bal_market}\\
X^{\mathrm{in}}_{qht}
  &= X^{\mathrm{out}}_{qht},
&& \forall q\in\setQ,\,h\in\setH,\,t\in\setT,
\label{eq:bal_hub}
\end{alignat}
\end{subequations}
with boundary condition $s^{\mathrm{inv}}_{pi,0}=I^0_{pi}$ for all
$p\in\setP$, $i\in\setNv$.

%- - - - - - - - - - - - - - - - - - - - - - - - - - - - - - - -
\subsubsection{(C2)~~Reverse logistics and returns}
\label{sec:model:c2}

Return flows are governed by product-specific collection rates and lag
parameters. The volume of recoverable item $k$ collected at market $r$
in period $t$ is bounded above and below by fractions of lagged
demand (with $d_{rp,t-\theta_p}\defeq 0$ when $t-\theta_p<1$). Every
recovered unit entering a factory is consumed entirely by remanufacturing,
closing the material loop:
\begin{subequations}\label{eq:c2_reverse}
\begin{alignat}{2}
\sum_{p\in\setP}\gamma_{kp}\,\underline{\eta}_{rpt}\,d_{rp,t-\theta_p}
  \le X^{\mathrm{out}}_{krt}
  &\le \sum_{p\in\setP}\gamma_{kp}\,\overline{\eta}_{rpt}\,d_{rp,t-\theta_p},
&& \forall k\in\setK,\,r\in\setR,\,t\in\setT,
\label{eq:return_bounds}\\
X^{\mathrm{in}}_{kft}
  &= \sum_{(p,\ell)\in\Omega^{\mathrm{r}}}
       \beta_{kp}\,q^{\mathrm{r}}_{pf\ell t},
&& \forall k\in\setK,\,f\in\setF,\,t\in\setT.
\label{eq:return_factory}
\end{alignat}
\end{subequations}

%- - - - - - - - - - - - - - - - - - - - - - - - - - - - - - - -
\subsubsection{(C3)~~Raw material supply}
\label{sec:model:c3}

A supplier can ship raw material $u$ in period $t$ only if it is both
open and activated via the binary $\sigma_{sut}$. The shipped quantity
is then bounded within the supplier's per-period capacity, and the
total raw-material inflow at each factory in period $t$ must exactly
satisfy production stoichiometry:
\begin{subequations}\label{eq:c3_supply}
\begin{alignat}{2}
\underline{S}_{su}\,\sigma_{sut}
  \le X^{\mathrm{out}}_{ust}
  &\le \overline{S}_{su}\,\sigma_{sut},
&& \forall s\in\setS,\,u\in\setU,\,t\in\setT,
\label{eq:supply_bounds}\\
\sigma_{sut}
  &\le y_s,
&& \forall s\in\setS,\,u\in\setU,\,t\in\setT,
\label{eq:supplier_link}\\
X^{\mathrm{in}}_{uft}
  &= \sum_{(p,g)\in\Omega^{\mathrm{p}}}
       \alpha_{upg}\,q^{\mathrm{p}}_{pfgt},
&& \forall u\in\setU,\,f\in\setF,\,t\in\setT.
\label{eq:raw_consumption}
\end{alignat}
\end{subequations}

%- - - - - - - - - - - - - - - - - - - - - - - - - - - - - - - -
\subsubsection{(C4)~~Facility opening and area capacity}
\label{sec:model:c4}

The floor area installed at each factory and warehouse is a continuous
decision that is active only when the facility is open. The installed
area is bounded by facility-specific limits conditional on opening, and
the used floor area $YCT_{it}$ (\Cref{eq:YCT}) -- combining handling
space for period throughput and, at warehouses, storage space for
end-of-period inventory -- must not exceed the installed area:
\begin{subequations}\label{eq:c4_capacity}
\begin{alignat}{2}
\underline{A}_i\,y_i
  \le YC_i
  &\le \overline{A}_i\,y_i,
&& \forall i\in\setF\cup\setW,
\label{eq:area_bounds}\\
YCT_{it}
  &\le YC_i,
&& \forall i\in\setF\cup\setW,\,t\in\setT.
\label{eq:capacity_use}
\end{alignat}
\end{subequations}

\begin{remark}[Flow-direction adaptation of the floor-area expression]
\label{rem:floor_area_flow_direction}
\citet{mota2018}'s general facility-area expression (their Eq.~16)
combines a flow-based handling term and a stock-based storage term,
$\sum_{m,a,j} apur_m X_{majit} + \sum_m apu_m S_{mit}$, where $X$
denotes flow \emph{into} node $i$ in their generic indexing
convention. This project adapts the flow-based term to its own
flow-direction convention rather than copying $X_{majit}$ verbatim.
At a factory, $X^{\mathrm{in}}_{pft}\equiv 0$ structurally under the
greenfield network-design convention of \Cref{rem:valid_ineq} --- no
arc routes finished product \emph{into} a factory --- so a literal
inbound-flow term would always vanish and defeat the purpose of
modelling factory handling space at all. Equation~\eqref{eq:YCT}
therefore uses the economically meaningful throughput quantity at a
factory, the outbound flow $X^{\mathrm{out}}_{pft}$, which by
\Cref{rem:valid_ineq} equals $q^{\mathrm{p}}_{pfgt}+q^{\mathrm{r}}_{pf\ell t}$
under that convention but remains well defined as a direct flow
quantity even if it did not. At a warehouse, both
$X^{\mathrm{in}}_{pwt}$ and $X^{\mathrm{out}}_{pwt}$ are generally
nonzero and differ only by the stock change~\eqref{eq:bal_warehouse},
so either is a defensible throughput proxy; \eqref{eq:YCT} retains the
inbound flow $X^{\mathrm{in}}_{pwt}$ as the more literal mirror of
Mota's own indexing convention, representing the receiving/staging
space needed for incoming goods, kept distinct from the
already-present stock term.
\end{remark}

%- - - - - - - - - - - - - - - - - - - - - - - - - - - - - - - -
\subsubsection{(C5)~~Inventory bounds}
\label{sec:model:c5}

Physical capacity and handling requirements impose per-product bounds
on end-of-period inventory. These bounds are active only when the
facility is open, linking the location and inventory decisions:
\begin{subequations}\label{eq:c5_inventory}
\begin{alignat}{2}
\underline{I}_{pi}\,y_i \le s^{\mathrm{inv}}_{pit}
  &\le \overline{I}_{pi}\,y_i,
&& \forall p\in\setP,\,i\in\setNv,\,t\in\setT.
\label{eq:inv_bounds}
\end{alignat}
\end{subequations}

%- - - - - - - - - - - - - - - - - - - - - - - - - - - - - - - -
\subsubsection{(C6)~~Technology selection and operation}
\label{sec:model:c6}

Production and remanufacturing technologies are selected at the factory
level and shared across all compatible products in set $\Omega^{\mathrm{p}}$
or $\Omega^{\mathrm{r}}$. Aggregate throughput per installed technology is
bounded in each period; remanufacturing operation in any period requires
prior installation; all technology installations require the factory to
be open; and at most one technology per product may be installed at each
factory:
\begin{subequations}\label{eq:c6_techcap}
\begin{alignat}{2}
\underline{Q}^{\mathrm{p}}_g\,\lambda^{\mathrm{p}}_{fg}
  \le \sum_{p:(p,g)\in\Omega^{\mathrm{p}}} q^{\mathrm{p}}_{pfgt}
  &\le \overline{Q}^{\mathrm{p}}_g\,\lambda^{\mathrm{p}}_{fg},
&& \forall f\in\setF,\,g\in\setG,\,t\in\setT,
\label{eq:prod_cap}\\
\underline{Q}^{\mathrm{r}}_{\ell}\,a_{f\ell t}
  \le \sum_{p:(p,\ell)\in\Omega^{\mathrm{r}}} q^{\mathrm{r}}_{pf\ell t}
  &\le \overline{Q}^{\mathrm{r}}_{\ell}\,a_{f\ell t},
&& \forall f\in\setF,\,\ell\in\setL,\,t\in\setT,
\label{eq:rem_cap}\\
a_{f\ell t}
  &\le \lambda^{\mathrm{r}}_{f\ell},
&& \forall f\in\setF,\,\ell\in\setL,\,t\in\setT,
\label{eq:rem_operation_link}\\
\lambda^{\mathrm{p}}_{fg},\;\lambda^{\mathrm{r}}_{f\ell}
  &\le y_f,
&& \forall f\in\setF,\,g\in\setG,\,\ell\in\setL,
\label{eq:tech_link}\\
\sum_{g:(p,g)\in\Omega^{\mathrm{p}}}\lambda^{\mathrm{p}}_{fg}
  &\le 1,
&& \forall p\in\setP,\,f\in\setF,
\label{eq:tech_exclusivity_prod}\\
\sum_{\ell:(p,\ell)\in\Omega^{\mathrm{r}}}\lambda^{\mathrm{r}}_{f\ell}
  &\le 1,
&& \forall p\in\setP,\,f\in\setF.
\label{eq:tech_exclusivity_rem}
\end{alignat}
\end{subequations}

%- - - - - - - - - - - - - - - - - - - - - - - - - - - - - - - -
\subsubsection{(C7)~~Transportation mode selection and shipments}
\label{sec:model:c7}

Transport services are contracted on each arc per period through the
binary mode-activation variable $\zeta_{ijmt}$. For owned road modes,
the total shipment weight must be consistent with the integer number
of trips $Q_{ijmt}$. Outsourced non-road modes are activated directly
by transported weight up to a contracted capacity, and at most one mode
may operate on any arc in any period:
\begin{subequations}\label{eq:c7_transport}
\begin{alignat}{2}
\underline{W}_m\,Q_{ijmt}
  \le \sum_{q:(q,i,j)\in\Phi}\omega_q\,x_{qijmt}
  &\le \overline{W}_m\,Q_{ijmt},
&& \forall (i,j)\in\setA,\,m\in\setM_{ij}\cap\setMrd,\,t\in\setT,
\label{eq:weight_trip}\\
\zeta_{ijmt} \le Q_{ijmt}
  &\le M^{\mathrm{trip}}_m\,\zeta_{ijmt},
&& \forall (i,j)\in\setA,\,m\in\setM_{ij}\cap\setMrd,\,t\in\setT,
\label{eq:trip_link}\\
\sum_{q:(q,i,j)\in\Phi}\omega_q\,x_{qijmt}
  &\le \overline{B}_m\,\zeta_{ijmt},
&& \forall (i,j)\in\setA,\,m\in\setM_{ij}\cap\setMout,\,t\in\setT,
\label{eq:contracted_active_cap}\\
\sum_{m\in\setM_{ij}}\zeta_{ijmt}
  &\le 1,
&& \forall (i,j)\in\setA,\,t\in\setT.
\label{eq:single_mode}
\end{alignat}
\end{subequations}

$M^{\mathrm{trip}}_m$ is stated above as a single constant per mode for
notational economy; the implementation derives a tightened,
arc- and node-specific bound instead (\Cref{app:impl_notes:bigm}).

%- - - - - - - - - - - - - - - - - - - - - - - - - - - - - - - -
\subsubsection{(C8)~~Fleet management and node activation}
\label{sec:model:c8}

Owned road vehicles must be physically stationed at their dispatching
node. A transport mode is usable only if both arc endpoints are active.
Outsourced modes require a period-level contract. The fleet at each
node is limited by driving-hour availability and by a total capital
budget. Finally, every open facility must carry positive flow to
prevent degenerate open-but-unused nodes:
\begin{subequations}\label{eq:c8_fleet}
\begin{alignat}{2}
\zeta_{ijmt} \le \bar{y}_{i},\quad
\zeta_{ijmt} &\le \bar{y}_{j},
&& \forall (i,j)\in\setA,\,m\in\setM_{ij},\,t\in\setT,
\label{eq:endpoint_act}\\
\zeta_{ijmt}
  &\le \kappa_{mt},
&& \forall (i,j)\in\setA,\,m\in\setM_{ij}\cap\setMout,\,t\in\setT,
\label{eq:contracted_cap}\\
\sum_{\substack{j:(i,j)\in\setA\\m\in\setM_{ij}}}
  \frac{r_m\,\delta_{ij}}{v_m}\,Q_{ijmt}
  &\le H_m\,K_{im},
&& \forall i\in\setN,\,m\in\setMrd,\,t\in\setT,
\label{eq:vehicle_hours}\\
K_{im} &\le M^{\mathrm{trip}}_m\,\bar{y}_i,
&& \forall i\in\setN,\,m\in\setMrd,
\label{eq:vehicle_act}\\
K_{im} &\le M^{\mathrm{trip}}_m
  \sum_{\substack{j:(i,j)\in\setA\\m\in\setM_{ij}}}
  \sum_{t\in\setT} Q_{ijmt},
&& \forall i\in\setN,\,m\in\setMrd,
\label{eq:vehicle_flow_link}\\
\sum_{i\in\setN}\sum_{m\in\setMrd}
  F^{\mathrm{m}}_m\,K_{im} &\le \overline{V},
\label{eq:fleet_budget}\\
\sum_{q\in\setQ}\sum_{m\in\setM}\sum_{t\in\setT}
  \sum_{\substack{(i,j)\in\setA\\m\in\setM_{ij}}}
    x_{qijmt}
  &\ge y_i,
&& \forall i\in\setNo.
\label{eq:min_flow}
\end{alignat}
\end{subequations}

Equation~\eqref{eq:vehicle_hours} presumes a positive mode speed
$v_m$ and a positive driving-hour allowance $H_m$; the
implementation's fallback for degenerate (non-positive) values is
given in \Cref{app:impl_notes:degenerate}.

%- - - - - - - - - - - - - - - - - - - - - - - - - - - - - - - -
\subsubsection{(C9)~~Intermodal hub routing}
\label{sec:model:c9}

Intermodal terminals act as mode-switching cross-docking points. At
each hub of type $\nu\in\{\mathrm{a},\mathrm{s}\}$, the volume arriving
by non-$\nu$ modes from non-$\nu$ sources must equal the volume
departing to same-type hubs by $\nu$-modes, and vice versa. Let
$\overline{\setM}^{\nu}_{ij}\defeq\setM_{ij}\cap\setM^{\nu}$ and
$\setM^{-\nu}_{ij}\defeq\setM_{ij}\setminus\setM^{\nu}$:
\begin{subequations}\label{eq:c9_hub}
\begin{alignat}{2}
\sum_{\substack{i:(i,h)\in\setA_q\\i\notin\setH^{\nu}}}
  \sum_{m\in\setM^{-\nu}_{ih}} x_{qihmt}
  &=
\sum_{\substack{j:(h,j)\in\setA_q\\j\in\setH^{\nu}}}
  \sum_{m\in\overline{\setM}^{\nu}_{hj}} x_{qhjmt},
&& \forall q\in\setQ,\,h\in\setH^{\nu},\,\nu,\,t\in\setT,
\label{eq:hub_ingress}\\
\sum_{\substack{i:(i,h)\in\setA_q\\i\in\setH^{\nu}}}
  \sum_{m\in\overline{\setM}^{\nu}_{ih}} x_{qihmt}
  &=
\sum_{\substack{j:(h,j)\in\setA_q\\j\notin\setH^{\nu}}}
  \sum_{m\in\setM^{-\nu}_{hj}} x_{qhjmt},
&& \forall q\in\setQ,\,h\in\setH^{\nu},\,\nu,\,t\in\setT.
\label{eq:hub_egress}
\end{alignat}
\end{subequations}

%- - - - - - - - - - - - - - - - - - - - - - - - - - - - - - - -
\subsubsection{(C10)~~Valid inequalities}
\label{sec:model:c10}

Every open factory must install at least one production or
remanufacturing technology,
\begin{equation}
y_f \le \sum_{g\in\setG}\lambda^{\mathrm{p}}_{fg}
       + \sum_{\ell\in\setL}\lambda^{\mathrm{r}}_{f\ell},
\qquad \forall f\in\setF,
\label{eq:vi_tech_demand}
\end{equation}
which is implied whenever the flow topology never routes finished
product \emph{into} a factory --- only the outbound arcs $\setA_{fw}$
exist for $p\in\setP$ (the standard convention for a greenfield
network-design instance, and the one used throughout
\Cref{sec:experiments}). A second candidate inequality, requiring
aggregate installed capacity to cover aggregate demand in every
period, is \emph{not} valid under inventory banking and is
deliberately excluded; \Cref{app:impl_notes:capdemand} gives the
equation and the counterexample.

\begin{remark}[Valid inequalities]\label{rem:valid_ineq}
Under the greenfield convention, $X^{\mathrm{in}}_{pft}=0$
structurally, so the flow-balance equation~\eqref{eq:bal_factory}
collapses to
$q^{\mathrm{p}}_{pfgt}+q^{\mathrm{r}}_{pf\ell t}=X^{\mathrm{out}}_{pft}$.
A factory installing neither technology therefore has
$X^{\mathrm{out}}_{pft}=0$ directly, and the raw-material and
recovered-item inflows that feed production and remanufacturing are
likewise pinned to zero by the equalities~\eqref{eq:raw_consumption}
and~\eqref{eq:return_factory}. Such a factory carries zero flow of
every item type, contradicting \eqref{eq:min_flow} if it is open. Note
that a single production technology alone --- the original, narrower
form of this inequality --- is \emph{not} implied: a factory engaged
purely in remanufacturing satisfies \eqref{eq:min_flow} via
$q^{\mathrm{r}}$ alone, without ever installing $\lambda^{\mathrm{p}}$.
Constraint~\eqref{eq:vi_tech_demand} is included in the implementation
because it is not redundant in the \LP relaxation; it tightens the
linear-programming bound and reduces the enumeration effort when the
\DBP solves the fixed-tax follower subproblems in
\Cref{sec:methodology}.
\end{remark}

%- - - - - - - - - - - - - - - - - - - - - - - - - - - - - - - -
\subsubsection{Domain constraints}
\label{sec:model:domains}

\begin{subequations}\label{eq:domains}
\begin{alignat}{1}
& y_i,\;\lambda^{\mathrm{p}}_{fg},\;\lambda^{\mathrm{r}}_{f\ell},\;
  a_{f\ell t},\;\sigma_{sut},\;\zeta_{ijmt},\;\kappa_{mt}
  \in \bin,
\label{eq:dom_bin}\\
& Q_{ijmt},\;K_{im} \in \Zplus,
\label{eq:dom_int}\\
& YC_i,\;q^{\mathrm{p}}_{pfgt},\;q^{\mathrm{r}}_{pf\ell t},\;
  s^{\mathrm{inv}}_{pit},\;x_{qijmt}
  \in \Rplus.
\label{eq:dom_cont}
\end{alignat}
\end{subequations}

Constraints~\eqref{eq:c1_balance}--\eqref{eq:domains} define the
follower's feasible region $\setX$. Together with the auxiliary
expressions~\eqref{eq:flows}--\eqref{eq:YCT}, the objectives
\eqref{eq:cost}--\eqref{eq:emissions}, the follower
programme~\eqref{eq:follower}, and the bilevel
programme~\eqref{eq:leader}, they constitute the \BLMIP solved in
\Cref{sec:methodology}.

%-----------------------------------------------------------------
% Backward-compatibility label for any cross-references to the old
% Follower's constraints subsection heading.
\label{sec:model:follower_constr}
